% !TeX program = xelatex
\documentclass{article}
\usepackage{kotex}
\usepackage[a4paper,left=3cm, right=3cm, top=2cm, bottom=2cm]{geometry}
\usepackage{amsmath}
\usepackage{graphicx}
\usepackage{caption}
\usepackage{setspace}
\usepackage{xcolor}
\usepackage{titlesec}
\usepackage{amssymb}
\usepackage{tcolorbox}
\usepackage{wrapfig}
\usepackage{amsthm}
\usepackage{multirow}
\usepackage{array}
\usepackage{tikz}
\usepackage{longtable}
\usepackage{pgfplots}
\pgfplotsset{compat=1.18}

\title{1.5 Elementary Matrices and a Method for Finding $A^{-1}$}
\date{}
\author{}
\setstretch{1.3}

\titleformat{name=\section, numberless}
  {\normalfont\large\bfseries\color{blue}}
  {}
  {0pt}
  {}
\geometry{a4paper, margin=1in}

% Red subheading style (matches the textbook's red section headers)
\newcommand{\redheading}[1]{{\vspace{2mm}\noindent\Large\bfseries\color{red!55!black} #1\par}\vspace{2mm}}

\newtheoremstyle{mystyle}
  {}{}{\itshape}{}{\bfseries}{.}{.5em}{}
\theoremstyle{mystyle}

\begin{document}
\maketitle

{\itshape\color{blue!45!black}
This section develops an algorithm for finding the inverse of a matrix, and along the way establishes further properties of invertible matrices.\par}

\redheading{Elementary Matrices}

Recall the three elementary row operations on a matrix $A$ (Section 1.1): multiply a row by a nonzero constant $c$; interchange two rows; add $c$ times one row to another. Each has an \textbf{inverse operation} that undoes it: multiply the same row by $1/c$; interchange the same two rows back; if $c$ times row $r_i$ was added to row $r_j$, add $-c$ times $r_i$ to $r_j$. Hence if $B$ results from $A$ by a sequence of row operations, some second sequence restores $A$ from $B$.

\begin{tcolorbox}[colback=teal!4, colframe=teal!65!black, title=Definition 1, fonttitle=\bfseries, boxrule=0.5mm, arc=3mm]
Matrices $A,B$ are \textbf{row equivalent} if either (hence each) can be obtained from the other by a sequence of elementary row operations.
\end{tcolorbox}

Next we show how matrix multiplication can perform a row operation.

\begin{tcolorbox}[colback=teal!4, colframe=teal!65!black, title=Definition 2, fonttitle=\bfseries, boxrule=0.5mm, arc=3mm]
A matrix $E$ is an \textbf{elementary matrix} if it results from $I$ by a \emph{single} elementary row operation.
\end{tcolorbox}

\subsection*{EXAMPLE 1 | Elementary Matrices and Row Operations}
\[
\begin{bmatrix}1&0\\0&-3\end{bmatrix},\quad
\begin{bmatrix}1&0&0&0\\0&0&0&1\\0&0&1&0\\0&1&0&0\end{bmatrix},\quad
\begin{bmatrix}1&0&3\\0&1&0\\0&0&1\end{bmatrix},\quad
\begin{bmatrix}1&0&0\\0&1&0\\0&0&1\end{bmatrix}
\]
result from, respectively: multiplying row 2 of $I_2$ by $-3$; interchanging rows 2, 4 of $I_4$; adding 3 times row 3 of $I_3$ to row 1; multiplying row 1 of $I_3$ by 1.

\begin{tcolorbox}[colback=red!4, colframe=red!60!black, title=Theorem 1.5.1 -- Row Operations by Matrix Multiplication, boxrule=0.5mm, arc=3mm]
If elementary matrix $E$ results from performing a row operation on $I_m$, and $A$ is $m\times n$, then $EA$ is the matrix that results from performing that same row operation on $A$.
\end{tcolorbox}

\subsection*{EXAMPLE 2 | Using Elementary Matrices}
For $A=\begin{bmatrix}1&0&2&3\\2&-1&3&6\\1&4&4&0\end{bmatrix}$ and $E=\begin{bmatrix}1&0&0\\0&1&0\\3&0&1\end{bmatrix}$ (which adds 3 times row 1 of $I_3$ to row 3):
\[
EA=\begin{bmatrix}1&0&2&3\\2&-1&3&6\\4&4&10&9\end{bmatrix}
\]
exactly the result of adding 3 times row 1 of $A$ to row 3.

Table 1 lists each row operation that produces $E$ from $I$, alongside its \textbf{inverse operation}, which restores $I$ from $E$.

\begin{center}
\begin{tabular}{|l|l|}
\hline
\textbf{Row Operation on $I$ That Produces $E$} & \textbf{Row Operation on $E$ That Reproduces $I$}\\
\hline
Multiply row $i$ by $c\ne0$ & Multiply row $i$ by $1/c$\\
\hline
Interchange rows $i$ and $j$ & Interchange rows $i$ and $j$\\
\hline
Add $c$ times row $i$ to row $j$ & Add $-c$ times row $i$ to row $j$\\
\hline
\end{tabular}
\end{center}

\subsection*{EXAMPLE 3 | Row Operations and Inverse Row Operations}
\[
\begin{bmatrix}1&0\\0&1\end{bmatrix}
\xrightarrow{\text{multiply row 2 by }7}
\begin{bmatrix}1&0\\0&7\end{bmatrix}
\xrightarrow{\text{multiply row 2 by }\tfrac17}
\begin{bmatrix}1&0\\0&1\end{bmatrix}
\]
\[
\begin{bmatrix}1&0\\0&1\end{bmatrix}
\xrightarrow{\text{interchange rows 1, 2}}
\begin{bmatrix}0&1\\1&0\end{bmatrix}
\xrightarrow{\text{interchange rows 1, 2}}
\begin{bmatrix}1&0\\0&1\end{bmatrix}
\]
\[
\begin{bmatrix}1&0\\0&1\end{bmatrix}
\xrightarrow{\text{add 5}\times\text{row 2 to row 1}}
\begin{bmatrix}1&5\\0&1\end{bmatrix}
\xrightarrow{\text{add }-5\times\text{row 2 to row 1}}
\begin{bmatrix}1&0\\0&1\end{bmatrix}
\]

\begin{tcolorbox}[colback=red!4, colframe=red!60!black, title=Theorem 1.5.2, boxrule=0.5mm, arc=3mm]
Every elementary matrix is invertible, and its inverse is also an elementary matrix.
\end{tcolorbox}
\textbf{Proof} Let $E$ result from a row operation on $I$, and let $E_0$ result from applying the \emph{inverse} operation to $I$. By Theorem 1.5.1 and the fact that inverse row operations cancel, $E_0E=I$ and $EE_0=I$; so $E_0=E^{-1}$. $\blacksquare$

\redheading{Equivalence Theorem}

Linear algebra's diverse ideas are deeply connected. The following theorem links invertibility, homogeneous systems, reduced row echelon form, and elementary matrices -- a first step in that direction (more will be added later).

\begin{tcolorbox}[colback=red!4, colframe=red!60!black, title=Theorem 1.5.3 -- Equivalent Statements, boxrule=0.5mm, arc=3mm]
For an $n\times n$ matrix $A$, the following are equivalent (all true or all false):\\
(a) $A$ is invertible.\quad \\
(b) $A\mathbf{x}=\mathbf{0}$ has only the trivial solution.\\
(c) The reduced row echelon form of $A$ is $I_n$.\quad \\
(d) $A$ is a product of elementary matrices.
\end{tcolorbox}
\textbf{Proof} We show $(a)\Rightarrow(b)\Rightarrow(c)\Rightarrow(d)\Rightarrow(a)$, which establishes that all four are equivalent.

$(a)\Rightarrow(b)$: If $A$ is invertible and $A\mathbf{x}_0=\mathbf{0}$, multiplying by $A^{-1}$ gives $\mathbf{x}_0=A^{-1}\mathbf{0}=\mathbf{0}$ -- only the trivial solution.

$(b)\Rightarrow(c)$: If $A\mathbf{x}=\mathbf{0}$ has only the trivial solution, Gauss--Jordan elimination reduces the augmented matrix $[A\mid\mathbf{0}]$ to $[I_n\mid\mathbf{0}]$ (each $x_i=0$); disregarding the zero last column, the reduced row echelon form of $A$ is $I_n$.

$(c)\Rightarrow(d)$: If $A$ reduces to $I_n$ by elementary row operations, then by Theorem 1.5.1 there are elementary matrices $E_1,\dots,E_k$ with $E_k\cdots E_1A=I_n$. Each $E_i$ is invertible (Theorem 1.5.2), so $A=E_1^{-1}\cdots E_k^{-1}$ -- a product of elementary matrices.

$(d)\Rightarrow(a)$: A product of elementary matrices is a product of invertible matrices (Theorem 1.5.2), hence invertible by Theorem 1.4.6. $\blacksquare$

\redheading{A Method for Inverting Matrices}

If $A$ is invertible, the elementary matrices reducing $A$ to $I_n$ satisfy $E_k\cdots E_1A=I_n$; multiplying on the right by $A^{-1}$ gives $A^{-1}=E_k\cdots E_1I_n$. That is, \emph{the same row operations that reduce $A$ to $I_n$ transform $I_n$ into $A^{-1}$.}

\begin{tcolorbox}[colback=blue!5, colframe=blue!60!black, title=Inversion Algorithm, boxrule=0.5mm, arc=3mm]
To find $A^{-1}$: find a sequence of elementary row operations reducing $A$ to $I$, and apply that same sequence to $I$ to obtain $A^{-1}$.
\end{tcolorbox}

In practice, adjoin $I$ to the right of $A$ to form $[A\mid I]$, then row-reduce until the left block is $I$; the right block is then $A^{-1}$, giving $[I\mid A^{-1}]$.

\subsection*{EXAMPLE 4 | Using Row Operations to Find $A^{-1}$}
Find the inverse of $A=\begin{bmatrix}1&2&3\\2&5&3\\1&0&8\end{bmatrix}$.

\textbf{SOLUTION:}
\[
\left[\begin{array}{ccc|ccc}1&2&3&1&0&0\\2&5&3&0&1&0\\1&0&8&0&0&1\end{array}\right]
\xrightarrow[-1R_1+R_3]{-2R_1+R_2}
\left[\begin{array}{ccc|ccc}1&2&3&1&0&0\\0&1&-3&-2&1&0\\0&-2&5&-1&0&1\end{array}\right]
\]
\[
\xrightarrow{2R_2+R_3}
\left[\begin{array}{ccc|ccc}1&2&3&1&0&0\\0&1&-3&-2&1&0\\0&0&-1&-5&2&1\end{array}\right]
\xrightarrow{-1\times R_3}
\left[\begin{array}{ccc|ccc}1&2&3&1&0&0\\0&1&-3&-2&1&0\\0&0&1&5&-2&-1\end{array}\right]
\]
\[
\xrightarrow[-3R_3+R_1]{3R_3+R_2}
\left[\begin{array}{ccc|ccc}1&2&0&-14&6&3\\0&1&0&13&-5&-3\\0&0&1&5&-2&-1\end{array}\right]
\xrightarrow{-2R_2+R_1}
\left[\begin{array}{ccc|ccc}1&0&0&-40&16&9\\0&1&0&13&-5&-3\\0&0&1&5&-2&-1\end{array}\right]
\]
Thus, $A^{-1}=\begin{bmatrix}-40&16&9\\13&-5&-3\\5&-2&-1\end{bmatrix}$.

If $A$ is not invertible, this will show up as a row of zeros appearing on the \emph{left} block at some stage -- by Theorem 1.5.3(a),(c), $A$ cannot then be invertible, and the computation may stop.

\subsection*{EXAMPLE 5 | Showing That a Matrix Is Not Invertible}
For $A=\begin{bmatrix}1&6&4\\2&4&-1\\-1&2&5\end{bmatrix}$:
\[
\left[\begin{array}{ccc|ccc}1&6&4&1&0&0\\2&4&-1&0&1&0\\-1&2&5&0&0&1\end{array}\right]
\xrightarrow[R_1+R_3]{-2R_1+R_2}
\left[\begin{array}{ccc|ccc}1&6&4&1&0&0\\0&-8&-9&-2&1&0\\0&8&9&1&0&1\end{array}\right]
\xrightarrow{R_2+R_3}
\left[\begin{array}{ccc|ccc}1&6&4&1&0&0\\0&-8&-9&-2&1&0\\0&0&0&-1&1&1\end{array}\right]
\]
A zero row appears on the left, so $A$ is not invertible.

\subsection*{EXAMPLE 6 | Analyzing Homogeneous Systems}
Determine whether each has nontrivial solutions:
(a) $x_1+2x_2+3x_3=0,\ 2x_1+5x_2+3x_3=0,\ x_1+8x_3=0$\qquad
(b) $x_1+6x_2+4x_3=0,\ 2x_1+4x_2-x_3=0,\ -x_1+2x_2+5x_3=0$

\textbf{SOLUTION:} By Theorem 1.5.3(a),(b), a homogeneous system has only the trivial solution iff its coefficient matrix is invertible. The coefficient matrix of (a) is the $A$ of Example 4, which is invertible, so (a) has only the trivial solution; the coefficient matrix of (b) is the $A$ of Example 5, which is not invertible, so (b) has nontrivial solutions.

\section*{Exercise Set 1.5}
\begin{enumerate}
    \item[1.] Determine whether the matrix is elementary.\\
    (a) $\begin{bmatrix}1&0\\-5&1\end{bmatrix}$\quad (b) $\begin{bmatrix}-5&1\\1&0\end{bmatrix}$\quad (c) $\begin{bmatrix}1&1&0\\0&0&1\\0&0&0\end{bmatrix}$\quad (d) $\begin{bmatrix}2&0&0&2\\0&1&0&0\\0&0&1&0\\0&0&0&1\end{bmatrix}$

    \item[3.] Find a row operation and its elementary matrix that restores the given elementary matrix to $I$.\\
    (a) $\begin{bmatrix}1&-3\\0&1\end{bmatrix}$\quad (b) $\begin{bmatrix}-7&0&0\\0&1&0\\0&0&1\end{bmatrix}$\quad (c) $\begin{bmatrix}1&0&0\\0&1&0\\-5&0&1\end{bmatrix}$\quad (d) $\begin{bmatrix}0&0&1&0\\0&1&0&0\\1&0&0&0\\0&0&0&1\end{bmatrix}$

    \item[5.] Identify the row operation corresponding to $E$ and verify that $EA$ results from applying it to $A$.\\
    (a) $E=\begin{bmatrix}0&1\\1&0\end{bmatrix}$, $A=\begin{bmatrix}-1&-2&5&-1\\3&-6&-6&-6\end{bmatrix}$\\
    (b) $E=\begin{bmatrix}1&0&0\\0&1&0\\0&-3&1\end{bmatrix}$, $A=\begin{bmatrix}2&-1&0&-4&-4\\1&-3&-1&5&3\\2&0&1&3&-1\end{bmatrix}$\\
    (c) $E=\begin{bmatrix}1&0&4\\0&1&0\\0&0&1\end{bmatrix}$, $A=\begin{bmatrix}1&4\\2&5\\3&6\end{bmatrix}$

    \item[7.] Using $A=\begin{bmatrix}3&4&1\\2&-7&-1\\8&1&5\end{bmatrix}$, $B=\begin{bmatrix}8&1&5\\2&-7&-1\\3&4&1\end{bmatrix}$, $C=\begin{bmatrix}3&4&1\\2&-7&-1\\2&-7&3\end{bmatrix}$, find an elementary matrix $E$ satisfying:\\
    (a) $EA=B$\quad (b) $EB=A$\quad (c) $EA=C$\quad (d) $EC=A$

    \item[9.] First use Theorem 1.4.5, then the inversion algorithm, to find $A^{-1}$ if it exists.\\
    (a) $A=\begin{bmatrix}1&4\\2&7\end{bmatrix}$\qquad (b) $A=\begin{bmatrix}2&-4\\-4&8\end{bmatrix}$

    \item[11.] Use the inversion algorithm to find the inverse, if it exists.\\
    (a) $\begin{bmatrix}1&2&3\\2&5&3\\1&0&8\end{bmatrix}$\qquad (b) $\begin{bmatrix}-1&3&-4\\2&4&1\\-4&2&-9\end{bmatrix}$

    \item[13.] Use the inversion algorithm to find the inverse, if it exists, of $\begin{bmatrix}1&0&1\\0&1&1\\1&1&0\end{bmatrix}$.

    \item[19.] Find the inverse of each $4\times4$ matrix, where $k_1,k_2,k_3,k_4,k$ are all nonzero.\\
    (a) $\begin{bmatrix}k_1&0&0&0\\0&k_2&0&0\\0&0&k_3&0\\0&0&0&k_4\end{bmatrix}$\qquad
    (b) $\begin{bmatrix}k&1&0&0\\0&1&0&0\\0&0&k&1\\0&0&0&1\end{bmatrix}$

    \item[21.] Find all values of $c$, if any, for which the matrix is invertible.
    \[
    \begin{bmatrix}c&c&c\\1&c&c\\1&1&c\end{bmatrix}
    \]

    \item[23.] Express the matrix and its inverse as products of elementary matrices.
    \[
    \begin{bmatrix}-3&1\\2&2\end{bmatrix}
    \]

    \item[27.] Show that $A=\begin{bmatrix}1&2&3\\1&4&1\\2&1&9\end{bmatrix}$ and $B=\begin{bmatrix}1&0&5\\0&2&-2\\1&1&4\end{bmatrix}$ are row equivalent by finding a sequence of elementary row operations that produces $B$ from $A$, and use it to find a matrix $C$ with $CA=B$.

    \item[29.] Show that if $A=\begin{bmatrix}1&0&0\\0&1&0\\a&b&c\end{bmatrix}$ is an elementary matrix, then at least one entry in the third row must be zero.

    \item[30.] Show that
    \[
    A=\begin{bmatrix}0&a&0&0&0\\b&0&c&0&0\\0&d&0&e&0\\0&0&f&0&g\\0&0&0&h&0\end{bmatrix}
    \]
    is not invertible for any values of the entries.

    \item[31.] \textbf{(Working with Proofs)} Prove that if $A,B$ are $m\times n$ matrices, then $A,B$ are row equivalent if and only if they have the same reduced row echelon form.

    \item[32.] \textbf{(Working with Proofs)} Prove that if $A$ is invertible and $B$ is row equivalent to $A$, then $B$ is also invertible.

    \item[33.] \textbf{(Working with Proofs)} Prove that if $B$ is obtained from $A$ by a sequence of elementary row operations, then a second sequence of elementary row operations, applied to $B$, recovers $A$.
\end{enumerate}

\end{document}
